\documentclass[letterpaper]{article}
\PassOptionsToPackage{unicode=true}{hyperref} % options for packages loaded elsewhere
\PassOptionsToPackage{hyphens}{url}
\usepackage{graphicx} % Required for inserting images


\def\tightlist{}


\graphicspath{ {../images/} }
\usepackage[margin=1in]{geometry}
\usepackage{tikz}
\usepackage[dvipsnames]{xcolor}
\usetikzlibrary{calc}




\usepackage[utf8]{inputenc}
\usepackage[T1]{fontenc}
\usepackage{lmodern}
\usepackage[english]{babel}

\usepackage{longtable}
\usepackage{booktabs}

% Useful packages
\usepackage{hyperref}
\usepackage{amsmath, amssymb}
\usepackage{fancyhdr}
\usepackage{setspace}
\usepackage{tocloft}
\usepackage{titlesec}


\usepackage{unicode-math}



% Code highlighting (Pandoc uses listings or minted, but listings is safer without shell escape)
\usepackage{listings}
\lstset{
  basicstyle=\ttfamily\small,
  breaklines=true,
  frame=single,
  backgroundcolor=\color[gray]{0.95}
}

\usepackage{setspace}


%Modifiable Commands; defaults set to none
\newcommand{\titleDOC}{}
\newcommand{\authorDOC}{Daniel Sabalakov}
\newcommand{\dateDOC}{\today}
\newcommand{\classDOC}{}
\newcommand{\emailDOC}{danielsabalakov@gmail.com}
% DEFAULT QUOTE IS BY ISSAC NEWTON. IF YOU DON'T WANT IT, DO: quote: @none
\newcommand{\quoteDOC}{``\textit{Truth is ever to be found in the simplicity, and not in the multiplicity and confusion of things}'' - Issac Newton}

% Header/Footer
\pagestyle{fancy}
\fancyhf{}
\rhead{\thepage}
\lhead{\leftmark}
\fancyhead[R]{\thepage} % Place page number in top right corner




% %Title Arguments
% \title{\TITLEOFDOCUMENT}
% \author{\AUTHOROFDOCUMENT}
% \date{\today}


%
% ANYTHING BELOW HERE CAN GET PROCEDURALLY GENERATED
%
% BEGINNING OF DOCUMENT
%
%
%
\begin{document}

%titlepage








\renewcommand{\quoteDOC}{``\textit{May the Force Be With You}'' \\ \textemdash{} General Jan Dodonna}

\renewcommand{\classDOC}{DE Introduction to Probability and Statistics, Portsmouth High School}

\renewcommand{\titleDOC}{Q1 Grade Analysis}

\renewcommand{\authorDOC}{Daniel Sabalakov}

\input{latex_templates/basetitle2.tex}

\section*{Overview}\label{overview}
\addcontentsline{toc}{section}{Overview}

Grades in Great Bay Stat. are categorized in three ways: Formative, Summative, Final. Each are weighted 20\%, 70\%, and 10\%, respectively; meaning that, technically, only the Summative Grade is required to pass the course. All grades analyzed will not include individual weights.

Let\textquotesingle s analyse Daniel Sabalakov\textquotesingle s, a student at Great Bay Stat., grades.

\section*{Data set}\label{data-set}
\addcontentsline{toc}{section}{Data set}

\subsection*{Formative}\label{formative}
\addcontentsline{toc}{subsection}{Formative}

Frequency Table for Formative:

\begin{longtable}[]{@{}llllll@{}}
\toprule\noalign{}
x & 92 & 93 & 95 & 100 & 101 \\
\midrule\noalign{}
\endhead
\bottomrule\noalign{}
\endlastfoot
N(X=x) & 1 & 1 & 1 & 7 & 3 \\
\end{longtable}

Raw Data for Formative Grades:
\{100,101,100,100,93,92,101,100,101,100,100,95,100\}

Formative Bar Graph and Box Plot

\begin{figure}[htbp]
	 \centering
	 \includegraphics[width=0.6\textwidth]{C:/Users/score/Downloads/form_bp.png}
\end{figure}

 \pagebreak 

\subsubsection*{Formative Mean and Standard Deviation, In Context}\label{formative-mean-and-standard-deviation-in-context}
\addcontentsline{toc}{subsubsection}{Formative Mean and Standard Deviation, In Context}

The mean of the Formative category is 98.692, and the standard deviation is 3.146. The standard deviation of 3.146 points is the typical distance between Daniel Sabalakov\textquotesingle s average Formative score of 98.692 points and one of Daniel Sabalakov\textquotesingle s scores.

\subsubsection*{Formative Median and IQR, In Context}\label{formative-median-and-iqr-in-context}
\addcontentsline{toc}{subsubsection}{Formative Median and IQR, In Context}

The median of the Formative category is 100 points, and the IQR of the Formative category is 0 points. (Q3-Q1=100-100=0) This means that the middle point of Daniel Sabalakov\textquotesingle s ordered points is a 100 and that the 25th percentile and the 75th percentile are 0 points apart.

\subsubsection*{Formative Outliers, In Context}\label{formative-outliers-in-context}
\addcontentsline{toc}{subsubsection}{Formative Outliers, In Context}

Using our two outlier rules of thumb, we can determine the left fence for each outlier.

\begin{enumerate}
\def\labelenumi{\arabic{enumi}.}
\tightlist
\item
  Mean \& Standard Deviation R.O.T: \(Outliers <,> \mu \pm 2\sigma\)

  \begin{itemize}
  \tightlist
  \item
    Outliers \textless{} \(98.692 - 2*3.146\) ==\textgreater{} Outliers \textless{} 92.4
  \item
    Outliers \textgreater{} \(98.692 + 2*3.146\) ==\textgreater{} Outliers \textgreater{} 104.984
  \end{itemize}
\item
  Median \& IQR R.O.T: \(Outliers <,> Q_{1,3} \pm 1.5(Q_3-Q_1)\)

  \begin{itemize}
  \tightlist
  \item
    Outliers \textless{} 100 - 1.5 * (100 - 100) ==\textgreater{} Outliers \textless{} 100
  \item
    Outliers \textgreater{} 100 + 1.5 * (100 - 100) ==\textgreater{} Outliers \textgreater{} 100
  \end{itemize}
\end{enumerate}

Using the Mean \& Standrard Deviations, scores under 92.4 are considered outliers for the data set. Daniel Sabalakov only had one score that is considered an outlier; the score of 92 is considered an outiler because 92 \textless{} 92.4

Using the Median \& IQR, scores under or over 100 are considered outliers. This yields a result of 6 outlier points. This means that points are very clustered around the 100 point mark since at least 50\% of data points in the Formative section are 100\textquotesingle s.

\subsection*{Summative}\label{summative}
\addcontentsline{toc}{subsection}{Summative}

Frequency Table for Summative Grades:

\begin{longtable}[]{@{}lllllll@{}}
\toprule\noalign{}
x & 75 & 86 & 95 & 96 & 100 & 101 \\
\midrule\noalign{}
\endhead
\bottomrule\noalign{}
\endlastfoot
N(X=x) & 1 & 1 & 4 & 1 & 4 & 1 \\
\end{longtable}

Raw Data for Summative Grades:
\{96,95,75,100,100,100,101,95,95,86,100,95\}

\begin{figure}[htbp]
	 \centering
	 \includegraphics[width=0.8\textwidth]{C:/Users/score/Downloads/sum_bg.png}
\end{figure}

\begin{figure}[htbp]
	 \centering
	 \includegraphics[width=0.8\textwidth]{C:/Users/score/Downloads/sum_bp.png}
\end{figure}

 \pagebreak 

\subsubsection*{Summative Mean and Standard Deviation, In Context}\label{summative-mean-and-standard-deviation-in-context}
\addcontentsline{toc}{subsubsection}{Summative Mean and Standard Deviation, In Context}

The mean of the Summative category is 94.833, and the standard deviation is 7.493. The standard deviation of 7.493 points is the typical distance between Daniel Sabalakov\textquotesingle s average Summative scores of 94.833 points and one of Daniel Sabalakov\textquotesingle s scores.

\subsubsection*{Summative Median and IQR, In Context}\label{summative-median-and-iqr-in-context}
\addcontentsline{toc}{subsubsection}{Summative Median and IQR, In Context}

The median of the Summative category is 95.5 points, and the IQR of the Summative category is 5 points (Q3-Q1=100-95=5) This means that the middle point of Daniel Sabalakov\textquotesingle s ordered points is a 95.5 and that the 25th percentile and the 75th percentile are 5 points apart.

\subsubsection*{Summative Outliers, In Context}\label{summative-outliers-in-context}
\addcontentsline{toc}{subsubsection}{Summative Outliers, In Context}

Using our two outlier rules of thumb, we can determine the left fence for each outlier.

\begin{enumerate}
\def\labelenumi{\arabic{enumi}.}
\tightlist
\item
  Mean \& Standard Deviation R.O.T: \(Outliers <,> \mu \pm 2\sigma\)

  \begin{itemize}
  \tightlist
  \item
    Outliers \textless{} \$94.833 - 2*7.493 \$ ==\textgreater{} Outliers \textless{} 79.846
  \item
    Outliers \textgreater{} \$94.833 + 2*7.493 \$ ==\textgreater{} Outliers \textgreater{} 109.819
  \end{itemize}
\item
  Median \& IQR R.O.T: \(Outliers <,> Q_{1,3} \pm 1.5(Q_3-Q_1)\)

  \begin{itemize}
  \tightlist
  \item
    Outliers \textless{} 95.5 - 1.5 * (5) ==\textgreater{} Outliers \textless{} 87.5
  \item
    Outliers \textgreater{} 95.5 + 1.5 * (5) ==\textgreater{} Outliers \textgreater{} 103
  \end{itemize}
\end{enumerate}

Using the Mean \& Standrard Deviations, scores under 79.846 are considered outliers for the Sumamtive data set. Daniel Sabalakov only had one score that is considered an outlier; the score of 75 is considered an outiler because 75 \textless{} 79.846.

Using the Median \& IQR, scores under a score of 87.5 or over a score of 103 are considered outliers. This yields a result of 2 outlier points: a score of 86 and a score of 75.

\subsection*{Final}\label{final}
\addcontentsline{toc}{subsection}{Final}

Daniel Sabalakov scored exceptionally well on the final, scoring a 100\% on each of the three portions of the final. This yields a median and a mean of a 100, and an IQR and Standard Deviation of 0.

\section*{Analysis and Conclusion}\label{analysis-and-conclusion}
\addcontentsline{toc}{section}{Analysis and Conclusion}

 \pagebreak 

\begin{figure}[htbp]
	 \centering
	 \includegraphics[width=0.8\textwidth]{C:/Users/score/Downloads/both_dp.png}
\end{figure}

A score of 75 in the Summative category, quite frankly, is an outlier beyond doubt. Calculating the z-score of a score of 75, we get: \(Z=\frac{y-\mu}{\sigma}\), ==\textgreater{} \(Z=\frac{75-94.833}{7.493}\) ==\textgreater{} \(Z=-2.669\). Using the Mean and Standard Deviation rule of thumb, any z-score under -2 is considered an outlier. If this is a normal model (which I will discuss in the next section), a Summative score of 75 or lower has a 0.38\% chance of occuring. Additionally, using the Median and IQR rule of thumb, any score under 87.5 is an outlier. While Daniel Sabalakov did score a 86, which would be considered an outlier, a score of 75 is even more of an outlier. Hence, a Summative score of 75 is undeniably an outlier from both the Median \& IQR perspective and a Mean and Standard Deviation perspective.

For data to be modeled normally, the following characteristics must be true.
Of the Summative distribution,

\begin{itemize}
\tightlist
\item
  Symmetry

  \begin{itemize}
  \tightlist
  \item
    The Summative distribution is \textbf{not} symmetric. An outlier of 75 points as well as an atypical 85 points make the curve left-skewed.
  \end{itemize}
\item
  Bell Shape (Unimodal):

  \begin{itemize}
  \tightlist
  \item
    The Summative distribution must be unimodal. The Summative Grade Distribution is \textbf{not} unimodal; rather, it is left-skewed.
  \end{itemize}
\item
  50\%-95\%-99\% rule

  \begin{itemize}
  \tightlist
  \item
    The Summative distribution does \textbf{not} follow the 50-95-99\% rule. A majority of the data is found within 1 standard deviation of the mean; 5/6th of the data is found within 1 standard deviation of the mean. 1/12 of the data (8.3\%) is found within 2 standard deviations, and 1/12 of the data (8.3\%) is an outlier. (Using the mean \& stdev R.O.T)
  \end{itemize}
\item
  Thin tails

  \begin{itemize}
  \tightlist
  \item
    The Summative distribution does \textbf{not} have thin tails on both asymptotes. On the right, the Summative Distribution is not thin-tailed.
  \end{itemize}
\end{itemize}

This means that the Summative distribution can not be modeled normally. However, this does not invalidate the point that a score of 75 is an outlier; a normal distribution only contributes a percentage.

Both the Summative and Formative score distributions are Left-skewed, meaning that the mean will be less than the median for both. If the Summative score of 75 (which is clearly an outlier) is removed from the Summative distribution, Daniel\textquotesingle s overall grade will be a 96.9\%. Removing the Summative score of 75 will additionally change the outlier rule of a typical distribution. The Summative distribution without the 75 has a mean of 96.63 with a standard deviation of 4.34 points. This means that, on average, a score in the modified distribution will vary by 4.34 points from the mean of 96.63 points.

\subsection*{Screenshot of Google Spreadsheet}\label{screenshot-of-google-spreadsheet}
\addcontentsline{toc}{subsection}{Screenshot of Google Spreadsheet}

\begin{figure}[htbp]
	 \centering
	 \includegraphics[width=1\textwidth]{C:/Users/score/Downloads/rstat1.png}
\end{figure}

\begin{figure}[htbp]
	 \centering
	 \includegraphics[width=1\textwidth]{C:/Users/score/Downloads/rstat2.png}
\end{figure}




\end{document}
